\begin{center}
    \textbf{Abstract}
\end{center}
%\phantomsection % required if using hyperref
%\addcontentsline{toc}{section}{\abstractname}
\doublespacing

Research on the impacts of mother tongue instruction (MTI) on labor market outcomes is limited. MTI improved salary employment in Ethiopia, without any effect on wage employment. The positive effect of MTI on salary employment is consistent with the prevailing ethno-linguistic federalism in the country and a suggestive evidence in this study that cognitive ability increased with MTI. Being from a region with a historically dominant language does not confer advantages in the labor market. MTI increases employment if it is complemented by appropriate institutional and policy measures.\\

\vspace{0in}
\noindent\textbf{Keywords:} Mother Tongue Instruction; Cognitive Ability; Employment;\\
\vspace{-.2in}\\
\noindent\textbf{JEL Codes:} I25, I28, J21, J24, J48\\

\bigskip

